\documentclass[letterpaper,journal]{IEEEtran}

\usepackage{amsmath,amsfonts}
\usepackage{array}
\usepackage{url}
\usepackage{graphicx}
\usepackage{cite}

\hyphenation{op-tical net-works semi-conduc-tor multi-modal token-izer token-izers}

\newcommand{\LAM}{large AI model}
\newcommand{\LAMs}{large AI models}
\newcommand{\GenAI}{generative AI}
\newcommand{\etal}{\emph{et al.}}
\newcommand{\etalspace}{\emph{et al.}\ }

\begin{document}

\title{From Generative to Embodied Intelligence: Network Evolution for Distributed Large-AI-Model Agents}

\author{Jihong Li, Bo Li,~\IEEEmembership{Student Member,~IEEE,}
Qingqing Wu,~\IEEEmembership{Senior Member,~IEEE,}
and Shunqing Zhang,~\IEEEmembership{Senior Member,~IEEE}%
\thanks{J. Li, B. Li, and S. Zhang are with the School of Communication and
Information Engineering, Shanghai University, Shanghai 200444, China.
Q. Wu is with the Department of Electronic Engineering, Shanghai Jiao Tong
University, Shanghai 200240, China. (Corresponding author: Shunqing Zhang.)}}

\markboth{IEEE Communications Surveys \& Tutorials,~Preprint}%
{Li \MakeLowercase{\textit{et al.}}: From Generative to Embodied Intelligence}

\maketitle

\begin{abstract}
Large AI models are evolving from generators that return digital content, to
agents that execute multi-step workflows, and further to embodied agents that
perceive, decide, and act in the physical world. This evolution changes not
only the volume but also the meaning, timing, and consequence of network
traffic. A generative request tolerates best-effort delivery if response
quality is preserved; an agentic workflow accumulates failure probability over
tool calls; and an embodied loop can become unsafe when multimodal observations
or action tokens are late, asynchronous, or distorted. This survey develops a
network-centric taxonomy across these three stages. It first aligns model
benchmarks with communication metrics, showing a progression from content
fidelity and token throughput, through end-to-end task reliability, to
closed-loop control cost, tail latency, and jitter. It then traces the
architecture evolution from cloud-centric inference to stateful edge
orchestration and hierarchical, multi-agent control. Enabling technologies are
organized into a token plane, a transport plane, and a coordination plane,
covering token compression and reuse, split and speculative inference,
finite-blocklength transmission, deterministic scheduling, semantic
communications, communication--control co-design, and dynamic model sharing.
A collaborative robotic-rescue case study demonstrates how the taxonomy
translates task and physical coupling into differentiated service-level
objectives. Finally, the paper identifies token recognition, tokenizer
compatibility, and token-native protocols as three prerequisites for networks
that connect intelligence rather than merely transport bits.
\end{abstract}

\begin{IEEEkeywords}
Agentic AI, embodied AI, edge intelligence, large AI models, multi-agent
systems, networked control, semantic communications, token communications,
ultra-reliable low-latency communications.
\end{IEEEkeywords}

\section{Introduction}
\label{sec:introduction}
\IEEEPARstart{L}{arge} AI models are moving beyond a single prompt--response
exchange. Generative models synthesize text, images, audio, and video;
agentic models decompose goals, invoke tools, retrieve memory, and cooperate
with other models; embodied models close the loop between multimodal sensing
and physical actuation. We use \emph{\LAM} (LAM) as an umbrella term for these
foundation-model-based systems, including large language models (LLMs),
multimodal large language models (MLLMs), world models, and
vision--language--action (VLA) models. The three stages are not mutually
exclusive model families. Rather, they represent increasingly stringent
interaction contracts between intelligence, a network, and an environment.

Mobile and edge networks are therefore changing roles. In the generative
stage, a network primarily transports prompts, context, and output tokens to a
remote accelerator. In the agentic stage, it also binds tools, model context,
persistent memory, and multiple expert models into a reliable transaction.
In the embodied stage, the network becomes part of the plant: it carries
camera, depth, point-cloud, inertial, joint-state, and action streams whose age
and synchronization affect the next physical state. The ongoing 3GPP study on
6G use cases explicitly includes AI-agent communication and network-supported
AI/ML training and inference, illustrating this shift from generic
connectivity toward AI, data, and computing service exposure
\cite{3gpp22870}.

This transition exposes a mismatch between conventional communication
objectives and application value. Maximizing throughput or minimizing average
packet delay does not specify whether a missing visual patch hides a grasp
point, whether a stale tool response invalidates the remaining plan, or
whether jitter destabilizes a fast control loop. Conversely, transmitting all
raw data with uniform reliability is often infeasible: an embodied platform
may contain multiple cameras and range sensors, tens of joints, and more than
one control timescale. The central design question is consequently:
\emph{which information must be delivered, to which model or actuator, by
what deadline, and with what task-level loss when resources are limited?}

\begin{figure*}[t]
\centering
\setlength{\tabcolsep}{2.5pt}
\renewcommand{\arraystretch}{1.15}
\begin{tabular}{c@{\hspace{4pt}$\Longrightarrow$\hspace{4pt}}c@{\hspace{4pt}$\Longrightarrow$\hspace{4pt}}c}
\fbox{\parbox[c][1.08in][c]{0.275\textwidth}{\centering
\textbf{Generative LAM}\\[2pt]
Prompt $\rightarrow$ content\\
one/few-shot interaction\\
quality, TTFT, token goodput\\[2pt]
\emph{dominant bottleneck: throughput}}}
&
\fbox{\parbox[c][1.08in][c]{0.275\textwidth}{\centering
\textbf{Agentic LAM}\\[2pt]
Goal $\rightarrow$ plan/tools/memory\\
multi-step and stateful\\
task success, tail job time\\[2pt]
\emph{dominant bottleneck: reliability}}}
&
\fbox{\parbox[c][1.08in][c]{0.275\textwidth}{\centering
\textbf{Embodied LAM}\\[2pt]
Sense $\rightarrow$ predict $\rightarrow$ act\\
continuous closed loop\\
control cost, jitter, safety\\[2pt]
\emph{bottleneck: delay--reliability}}}
\\[5pt]
\multicolumn{3}{c}{\small Increasing interaction frequency, source coupling,
physical consequence, and network--compute--control interdependence}
\end{tabular}
\caption{Evolution of LAM services and their dominant network operating
bottlenecks. Requirements accumulate rather than replace one another:
embodied systems still need throughput and semantic reliability, but must
add deterministic closed-loop behavior.}
\label{fig:evolution}
\end{figure*}

\subsection{Scope and Survey Method}
This paper is a structured technical survey at the intersection of four
literatures that are often treated separately: (i) model evaluation and
efficient inference, (ii) edge intelligence and semantic communications,
(iii) ultra-reliable low-latency communications (URLLC), and (iv) wireless
networked control. We prioritize peer-reviewed papers, archival preprints for
rapidly emerging VLA techniques, and active standards. Studies are compared
along a common chain:
\[
\text{model}\rightarrow\text{token traffic}\rightarrow
\text{network SLO}\rightarrow\text{physical task}.
\]
The goal is not to rank model leaderboards or to repeat a general edge-AI
survey. It is to reveal where a metric, architecture, or mechanism developed
for one stage fails when it is transferred to the next.

\subsection{Positioning Relative to Adjacent Surveys}
The topic is adjacent to, but broader than, several established survey areas.
Edge-LLM studies primarily ask where a model should execute and how model,
tensor, pipeline, or expert parallelism reduces inference cost. They usually
treat the model output as the terminal objective. Semantic-communication
studies replace bit fidelity with meaning or task fidelity, and recent token
communication makes foundation-model tokens the exchanged representation
\cite{semanticclassic,tokcom}. Integrated sensing and edge AI extends the
objective to perception tasks and systematically co-designs sensing,
communication, and inference \cite{isea}. URLLC and deterministic networking
provide reliability and tail-delay tools, while wireless networked control
directly optimizes stability, identification, or linear-quadratic cost.

Each area solves one part of the embodied problem, but their boundary
assumptions differ. Efficient LLM inference generally assumes that additional
latency reduces user experience but not plant stability. Semantic
communications commonly evaluates a reconstructed signal or a downstream
classifier, for which occasional errors do not change the next source state.
URLLC abstracts a packet source and deadline without asking which physical
state generated the packet. Networked-control theory closes the loop but
usually starts with an explicit low-dimensional dynamical model rather than a
VLA policy and multimodal token stream.

This survey uses model evolution to connect the areas. Its distinctive object
is the \emph{distributed LAM agent service}: model execution, network
transport, memory/tool interaction, and physical control are all parts of one
service graph. Table~\ref{tab:positioning} makes the boundary explicit. The
comparison is important because merely combining the mechanisms in adjacent
columns does not ensure end-to-end validity. For example, a token-pruning
method can preserve perception accuracy, and a URLLC link can deliver the
retained tokens reliably, while the closed loop still fails because the
pruned token marks a contact transition or the accelerator queue introduces
jitter.

\begin{table*}[t]
\caption{Position of This Survey Relative to Adjacent Research Areas}
\label{tab:positioning}
\centering
\footnotesize
\renewcommand{\arraystretch}{1.17}
\setlength{\tabcolsep}{3.1pt}
\begin{tabular}{|p{0.15\textwidth}|p{0.18\textwidth}|p{0.18\textwidth}|p{0.18\textwidth}|p{0.21\textwidth}|}
\hline
\textbf{Area} & \textbf{Terminal objective} & \textbf{Traffic abstraction}
& \textbf{Typical system boundary} & \textbf{What remains outside that
boundary} \\ \hline
Edge LLM inference &
Inference latency, energy, accuracy, or cost &
Requests, tensors, features, experts, and output tokens &
Device--edge--cloud model serving &
Tool transactions, control dynamics, task-dependent physical consequence
\\ \hline
Semantic/token communications &
Semantic reconstruction or downstream inference quality &
Meaning-bearing features or model tokens &
Semantic encoder, channel, and semantic decoder &
Long-horizon task dependency, actuator timing, and closed-loop stability
\\ \hline
Integrated sensing and edge AI &
End sensing/perception task performance &
Sensor signals and task-oriented features &
Sensing--communication--edge inference &
Agentic tools/memory and multi-rate action execution \\ \hline
URLLC/deterministic networking &
Packet error, hard deadline, delay tail, and jitter &
Short packets and queues with service classes &
Radio/access and sometimes edge-compute queue &
Source semantics, cross-modal dependence, and physical-state evolution
\\ \hline
Wireless networked control &
Stability, estimation/identification error, or control cost &
State estimates, measurements, and control packets &
Communication link plus an explicit dynamical plant/controller &
Foundation-model uncertainty, tokenization, and implicit learned dynamics
\\ \hline
\textbf{This survey} &
Content quality $\rightarrow$ task success $\rightarrow$ safe physical
control &
Tokens, transactions, multimodal observations, actions, and dependencies &
Model--memory/tool--network--compute--physical service graph &
Develops the interfaces and open problems needed to optimize the graph
end-to-end \\ \hline
\end{tabular}
\end{table*}

\subsection{Contributions}
The survey makes four contributions.

\begin{itemize}
\item It proposes a three-stage taxonomy that jointly classifies model
objectives, token sources, interaction patterns, benchmarks, failure
consequences, and network operating bottlenecks.
\item It maps AI-side evaluation to network-side quality of service (QoS).
The resulting metric ladder progresses from time to first token and content
quality, to transaction success and long-horizon tail latency, and finally to
task distortion, robust loop period, and control cost.
\item It organizes enabling technologies into token, transport, and
coordination planes, and explains their cross-layer coupling. This highlights
two gaps in current work: compression is commonly perception-oriented rather
than control-oriented, and deterministic transmission is commonly
source-independent rather than task-coupled.
\item It translates the taxonomy into an embodied multi-agent rescue case
study and derives research directions for token-aware service classification,
interoperability, and protocol design.
\end{itemize}

Section~\ref{sec:model_evolution} reviews model evolution and benchmarks.
Sections~\ref{sec:qos} and~\ref{sec:architecture} study metric and architecture
evolution, respectively. Section~\ref{sec:enablers} surveys enabling
technologies, Section~\ref{sec:case} presents the case study, and
Section~\ref{sec:future} discusses open challenges.

\section{From Generative to Embodied LAMs: Performance Metrics, Benchmarks, and Model Evolution}
\label{sec:model_evolution}

\subsection{Generative LAMs}
Generative LAMs learn a conditional distribution over token sequences.
Text is represented by subword tokens; images and video by spatial or
spatiotemporal patches or discrete latent codes; and audio by waveform,
spectral, or codec tokens. Transformer attention supplies a shared processing
interface, but it does not make these tokens equivalent from a networking
perspective. A text token may represent a word fragment, whereas the loss of a
visual token can remove localized geometry and the loss of an audio token can
corrupt a time interval.

Model-side evaluation in this stage is largely \emph{content-centric}.
Perplexity measures predictive fit, exact-match or multiple-choice accuracy
measures knowledge and reasoning, and BLEU/ROUGE-like scores compare generated
and reference text. MMLU spans 57 subjects and tests broad academic knowledge
\cite{mmlu}; BIG-bench uses more than 200 heterogeneous tasks to probe
capabilities and scaling behavior \cite{bigbench}. For open-ended and
multimodal generation, human preference, learned reward, semantic similarity,
factuality, and safety are required because surface-form overlap is
insufficient.

The corresponding traffic is asymmetric and elastic. Prompts and retrieved
context form a relatively bursty uplink; autoregressive output creates a
streaming downlink. A request can often tolerate a longer queue if the result
remains useful. Thus cloud serving, request batching, key--value cache
management, model placement, and output-token scheduling dominate. Edge
inference reduces backhaul and privacy exposure but introduces memory,
compute, and wireless partitioning constraints
\cite{beyondcloud,modeloffload}. The
fundamental network tradeoff is \emph{quality--latency--cost}: compression,
quantization, or lossy semantic delivery is acceptable when the generated
content meets the user's intent.

\subsection{Agentic LAMs}
An agentic LAM places a generative model inside a feedback loop with a planner,
tool interfaces, memory, verification, and sometimes other agents. The unit of
work is no longer one output sequence but a trajectory
\[
\tau=(o_0,a_0,o_1,a_1,\ldots,o_K),
\]
where observations $o_k$ may be web pages, application screenshots, database
results, code-execution logs, or messages from peer agents, and $a_k$ may call
an API or change external state. Reasoning techniques such as chain-of-thought
and tree/graph search, retrieval-augmented generation, model context
protocols, and multi-agent debate increase capability by introducing more
rounds and more state.

Consequently, evaluation shifts from answer similarity to
\emph{process-grounded task completion}. AgentBench, for example, evaluates
LLM agents across operating-system, database, knowledge-graph, gaming,
household, web-shopping, and web-browsing environments
\cite{agentbench}. Relevant metrics include task success rate, number of tool
calls, constraint violations, plan length, recovery from tool failure,
long-horizon consistency, and cost per completed task. An agent that produces
fluent intermediate text but executes the wrong transaction has failed.

Network reliability compounds over the trajectory. If step $k$ succeeds with
conditional probability $p_k$, a simplified serial workflow has
\begin{equation}
P_{\rm task}=\prod_{k=1}^{K}p_k,\qquad
P_{\rm fail}\leq\sum_{k=1}^{K}(1-p_k),
\label{eq:agent_reliability}
\end{equation}
where the equality assumes that each $p_k$ already conditions on prior
success. Even modest per-call error becomes material for large $K$.
Retries improve reliability but expand context, consume tokens, and create a
heavy tail in job completion time. Heterogeneous model collaboration and
adaptive repeat-query mechanisms can trade extra inference rounds for higher
task reliability \cite{garq}. Agentic networks must therefore expose
transaction identifiers, tool semantics, memory freshness, idempotence, and
provenance rather than treating every model call as an independent flow.

\subsection{Embodied LAMs}
Embodied LAMs ground language and vision in state estimation, world
prediction, planning, and action. Architectures include VLA models, latent
world models, bird's-eye-view transformers, and joint-embedding predictive
architectures. The observation $o_t$ can combine RGB/depth images, point
clouds, force/torque, inertial measurements, joint encoders, and messages from
other robots. The output $u_t$ can contain waypoints, discrete skills, joint
targets, or continuous control actions. Generalist humanoid models such as
GR00T N1 illustrate a trend toward shared representations across embodiments
and tasks \cite{groot}.

The benchmark moves into a physical or high-fidelity simulated loop.
Meta-World contains 50 robotic manipulation tasks and evaluates multi-task
generalization \cite{metaworld}; modern VLA evaluations additionally report
task success, trajectory efficiency, collision rate, action error, robustness
to scene perturbation, and sim-to-real degradation. World models are judged
by predictive consistency and planning quality, not merely reconstruction.
Because a small perception error is harmless in one state and catastrophic
near a contact transition, average token distortion alone is inadequate.

The control-oriented objective can be written as
\begin{equation}
J_{\rm c}=\mathbb{E}\!\left[\sum_{t=0}^{H}
\gamma^t\!\left(
\|x_t-x_t^\star\|_{\mathbf Q}^{2}
+\|u_t\|_{\mathbf R}^{2}
+\lambda_{\rm s}{\bf 1}\{\text{unsafe at }t\}
\right)\right],
\label{eq:control_cost}
\end{equation}
where $x_t$ is the physical state, $x_t^\star$ the task-dependent target,
$u_t$ the action, and ${\bf 1}\{\cdot\}$ is an indicator that penalizes safety
constraint violation. The important point is that communication changes
$J_{\rm c}$ through delayed, missing, or distorted observations and actions.

Embodied traffic also has structure unavailable in generic Internet flows.
Visual tokens often contain \emph{spatial} redundancy because only a subset
describes task-relevant objects. Action tokens often contain \emph{temporal}
redundancy during stable motion but become highly sensitive near transitions.
FlashVLA exploits both visual-token pruning and action reuse, reporting
substantial compute and latency savings with a small success-rate reduction
on LIBERO \cite{flashvla}. Yet temporal redundancy must not be confused with
permission to drop small controls: near an unstable equilibrium, small
corrective actions can be essential.

\begin{table*}[t]
\caption{Model and Traffic Evolution Across LAM Stages}
\label{tab:modeltaxonomy}
\centering
\footnotesize
\renewcommand{\arraystretch}{1.18}
\setlength{\tabcolsep}{3.2pt}
\begin{tabular}{|p{0.085\textwidth}|p{0.17\textwidth}|p{0.18\textwidth}|p{0.17\textwidth}|p{0.16\textwidth}|p{0.15\textwidth}|}
\hline
\textbf{Stage} & \textbf{Primary objective} & \textbf{Token/source structure}
& \textbf{Representative evaluation} & \textbf{Traffic pattern}
& \textbf{Failure consequence} \\ \hline
Generative &
Produce useful and faithful digital content &
Subwords, image/video patches, audio or codec tokens; statistical and
cross-modal redundancy &
Perplexity, accuracy, MMLU, BIG-bench, BLEU/ROUGE, preference, factuality &
Bursty prompt/context uplink and streaming output; one/few-shot; elastic
deadline &
Hallucinated or degraded content; often detectable, retryable, and
non-physical \\ \hline
Agentic &
Complete a goal through reasoning, tools, memory, and collaboration &
Generated tokens plus UI frames, logs, API objects, retrieved memory, and
inter-agent messages; workflow-induced redundancy &
Task success, tool correctness, constraint violations, calls/steps, recovery,
cost per completed job &
Stateful multi-round graph with fan-out, retries, persistent context, and
long-tail completion time &
Wrong external operation or cascading plan failure; compensation may be
expensive or impossible \\ \hline
Embodied &
Perceive, predict, and act safely in the physical world &
Multiview/multimodal perception, world-state and action tokens; spatial,
temporal, task, and physical coupling &
Task success, action/trajectory error, collision/safety, world-model error,
planning quality, sim-to-real gap, control cost &
Continuous, bidirectional, multisource loops at heterogeneous periods;
synchronization- and jitter-sensitive &
Unstable motion, equipment damage, or injury; late information can be worse
than missing information \\ \hline
\end{tabular}
\end{table*}

\subsection{Cross-Stage Synthesis}
Table~\ref{tab:modeltaxonomy} reveals three trends. First, evaluation moves
from \emph{relative content accuracy}, to \emph{practical process success}, to
\emph{physical closed-loop performance}. Second, interaction frequency and
coupling increase: a prompt is weakly coupled to the next request, an agent
trajectory is logically coupled across steps, and an embodied trajectory is
both logically and dynamically coupled through the environment. Third, the
consequence of failure changes from a correctable response to an external
state change and finally to a potentially irreversible physical event.

The increasing regularity of traffic creates an opportunity as well as a
burden. Generative tokens mainly expose statistical redundancy; agentic tokens
add workflow and protocol redundancy; embodied tokens add constraints induced
by geometry, kinematics, and task relations. A network that can recognize
these structures can avoid transporting semantically predictable tokens.
However, it must evaluate prediction errors in the task domain. This motivates
the QoS evolution in the next section.

\subsection{Benchmark--Network Co-Design}
A benchmark becomes a useful network workload only when its execution
semantics are preserved. Three mappings are required. First, the benchmark
unit must be declared. For generation it may be a prompt and response; for an
agent it is a complete trajectory including tool side effects; for a robot it
is an episode with initial state, perturbations, and termination criteria.
Reporting per-call latency for an agent or per-frame accuracy for a robot
silently changes the evaluation unit and can reverse conclusions.

Second, perturbations must be injected at the layer where a proposed mechanism
acts. A radio study should vary loss bursts, rate, queuing, handover, and clock
skew; a split-inference study should additionally vary feature quantization,
device stragglers, and accelerator queues. These perturbations must retain
their temporal structure. Independently dropping 1\% of packets is not
equivalent to a short outage during grasp contact, and randomly shuffling a
mean-delay trace destroys the jitter and burst correlations that a controller
experiences.

Third, the benchmark must propagate the perturbation to the terminal metric.
For a generative model this requires decoding the received representation and
scoring the output, not only measuring feature mean-square error. For an agent
it requires replaying the remaining task graph, because a delayed observation
can alter later tool calls. For an embodied agent it requires closing the
physics loop. DaDu-Corki's algorithm--architecture analysis exemplifies how
embodied-model inference latency must be studied together with robotic
execution rather than as an isolated transformer kernel \cite{dadu}.

The mapping also reveals a reproducibility hierarchy. A \emph{trace replay}
uses fixed model and network traces and is inexpensive, but cannot represent
feedback. A \emph{co-simulation} couples real model inference and network
emulation to simulated tools or physics. A \emph{hardware-in-the-loop} setup
adds real sensors, radio, or actuators and captures release-time and clock
effects. A \emph{field trial} tests distribution shift and safety procedures.
Claims should match the highest validated level. In particular, a simulation
gain in task success is evidence for an orchestration principle, not proof of
real-world safety.

Finally, benchmark scores should be reported as a surface rather than a
single operating point. The surface relates model scale, token budget,
wireless resources, compute load, and task quality. It exposes Pareto fronts
and phase transitions---for example, a stable range in which action reuse is
safe followed by a sharp loss near contact. Such a surface is the empirical
counterpart of the task-conditioned rate--distortion and robust-loop
formulations developed later in this paper.

\section{From Generative to Embodied Networks for LAMs: QoS Metric Evolution}
\label{sec:qos}

\subsection{QoS for Generative-LAM Networks}
Generative serving decomposes user-perceived delay into request admission,
prefill, and decode. The most useful operational metrics are time to first
token (TTFT), time per output token (TPOT), end-to-end response time, token
goodput, and deadline attainment. Token goodput counts tokens that satisfy a
service-level objective (SLO), rather than raw decoded tokens. It prevents a
server from appearing efficient by admitting too many requests and violating
their latency targets.

Quality and cost must accompany latency. Let $q_i$ be a task-specific response
quality, $T_i^{\rm f}$ the TTFT, $r_i^{\rm tok}$ the output rate, and $e_i$ the
energy or monetary cost for request $i$. A generic utility is
\begin{equation}
U_i^{\rm gen}=w_q q_i-w_f[T_i^{\rm f}-d_i^{\rm f}]^+
+w_r\log r_i^{\rm tok}-w_e e_i,
\label{eq:gen_utility}
\end{equation}
where $[z]^+=\max(z,0)$. The weights represent a service tier. Interactive
chat emphasizes TTFT and TPOT; offline generation emphasizes completion time
and energy; high-fidelity media emphasizes semantic quality and sustained
throughput. Byte throughput alone cannot represent these differences.

\subsection{QoS for Agentic-LAM Networks}
For an agentic service, each model response is an intermediate event. QoS must
be measured at the transaction boundary: end-to-end job completion time,
success probability, recovery probability, number of unnecessary calls,
context freshness, and resource cost. Percentiles matter because branching,
tool timeouts, and retries generate a long-tailed latency distribution.

The network should distinguish \emph{semantic reliability} from packet
reliability. A packet can be delivered correctly while a stale observation,
duplicated non-idempotent tool call, incompatible schema, or poisoned memory
causes task failure. Conversely, a lost low-value rationale token need not
trigger retransmission if a verified tool result is intact. A useful
agentic-SLO tuple is
\[
{\cal S}_{\rm ag}=(P_{\rm task},T_{95}^{\rm job},A_{\rm mem},
C_{\rm job},P_{\rm unsafe}),
\]
where $A_{\rm mem}$ is the age of task-relevant memory and $C_{\rm job}$
accounts for communication, inference, storage, and tool access. Network-aware
model-context platforms are beginning to make capability discovery and
network state visible to an agent \cite{netmcp}; a complete design must also
support admission, rollback, and audit across the full task graph.

\subsection{QoS for Embodied-LAM Networks}
An embodied loop adds sensing, computation, communication, and actuation:
\begin{equation}
T_t^{\rm loop}=T_t^{\rm sense}+T_t^{\rm enc}+T_t^{\rm ul}
+T_t^{\rm queue}+T_t^{\rm inf}+T_t^{\rm dl}+T_t^{\rm act}.
\label{eq:loop_delay}
\end{equation}
Optimizing one component without accounting for the others can be
counterproductive. For example, stronger visual compression may reduce
$T^{\rm ul}$ but increase $T^{\rm enc}$ and action error; offloading can
reduce $T^{\rm inf}$ but add queue and radio tails.

Average delay is insufficient. For a loop period $h_i$, tolerance
$\epsilon_i$, and violation probability $\delta_i$, a robust closed-loop SLO
can be expressed as
\begin{equation}
\Pr\!\left\{T_{i,t}^{\rm loop}\le h_i,\,
\left|T_{i,t}^{\rm loop}-\mathbb{E}T_{i,t}^{\rm loop}\right|
\le\epsilon_i h_i\right\}\ge 1-\delta_i.
\label{eq:robust_slo}
\end{equation}
The first condition controls deadline reliability; the second limits jitter.
For asynchronous multi-rate systems, relative age and skew among modalities
must also be bounded. A perfectly delivered camera frame timestamped before a
rapid contact change can be less useful than a lower-resolution but current
frame.

The proper distortion is task-conditional. For a retained token set
${\cal K}$, define
\begin{equation}
D_{\rm task}({\cal K}\mid g,x)
=J_{\rm c}({\cal K}\mid g,x)-J_{\rm c}({\cal K}_{\rm full}\mid g,x),
\label{eq:task_distortion}
\end{equation}
where $g$ is the goal and $x$ the current physical state. The same token can
have low distortion while navigating free space and high distortion near an
obstacle. This state dependence explains why static top-$k$ selection or
feature-reconstruction error does not directly optimize embodied performance.

Finally, the economic metric should be \emph{control utility per resource
cost}, not bits per joule alone:
\begin{equation}
\eta_{\rm EAI}=
\frac{J_{\rm c}^{\rm baseline}-J_{\rm c}^{\rm scheme}}
{\alpha B+\beta E+\chi F+\zeta S},
\label{eq:cost_effectiveness}
\end{equation}
where $B,E,F,S$ denote bandwidth, energy, computing cycles, and storage.
Equation~\eqref{eq:cost_effectiveness} discourages spending large resources on
tokens that do not improve the physical task.

\begin{table*}[t]
\caption{QoS Ladder and Representative Observability}
\label{tab:qos}
\centering
\footnotesize
\renewcommand{\arraystretch}{1.18}
\setlength{\tabcolsep}{3.2pt}
\begin{tabular}{|p{0.09\textwidth}|p{0.21\textwidth}|p{0.19\textwidth}|p{0.19\textwidth}|p{0.22\textwidth}|}
\hline
\textbf{Stage} & \textbf{Primary SLOs} & \textbf{Useful secondary metrics}
& \textbf{Insufficient legacy proxy} & \textbf{Required cross-layer
observability} \\ \hline
Generative &
TTFT, TPOT, completion deadline, semantic quality, token goodput &
Batch efficiency, cache hit ratio, energy/token, price/request &
Average throughput without response-quality or deadline qualification &
Prompt/output length, model and quantization, cache state, queue, radio rate
\\ \hline
Agentic &
Task success, $T_{95/99}^{\rm job}$, unsafe-operation probability, cost/job &
Tool-call count, retry depth, memory age, recovery and rollback time &
Per-call packet success or average model latency &
Task graph, dependency and idempotence, tool status, context version,
model confidence/provenance \\ \hline
Embodied &
Control cost, task success, loop deadline reliability, jitter, modality skew,
safety violation &
Action error, age/value of information, sim-to-real gap, energy/task,
cost-effectiveness &
Average one-way latency, raw PSNR, or independent per-flow reliability &
Goal and control phase, token modality, physical/task coupling, timestamps,
uncertainty, compute/radio/actuation state \\ \hline
\end{tabular}
\end{table*}

\subsection{Why the Metrics Cannot Be Collapsed into One KPI}
A single scalar ``AI QoS'' hides operationally important tradeoffs.
Generative quality is often evaluated after a response is complete; agentic
reliability is path-dependent; embodied value changes at the control
timescale. Network optimization should therefore maintain a vector SLO and
apply a task-specific scalarization only at admission or scheduling time.
Table~\ref{tab:qos} summarizes the necessary observability. The progression is
from flow state, to transaction state, to physical state. This is also the
progression in how much application information a network must obtain,
creating privacy and layering challenges discussed in
Section~\ref{sec:future}.

\subsection{An End-to-End Measurement Contract}
Cross-layer optimization is impossible without a common clock and causal
identifiers. We propose a measurement contract with four linked records:
\emph{information}, \emph{execution}, \emph{transport}, and \emph{outcome}.
An information record identifies a modality or token chunk, its generation
time, model/tokenizer version, uncertainty, dependency, and declared
deadline. An execution record captures model placement, queue entry/exit,
batch, cache status, and accelerator time. A transport record captures packet
size, path, coding/retransmission, radio service, and receiver timestamp. An
outcome record closes the chain with response quality, tool commit, action,
task state, or control loss.

The records need a shared \emph{causal identifier}, not merely synchronized
timestamps. A single observation can generate several candidate actions; an
agent retry can reuse an earlier observation; and a batched accelerator can
serve chunks from different tasks. Without parent--child identifiers, one
cannot distinguish whether a late action resulted from radio congestion,
batch waiting, a model escalation, or a tool retry. This matters both for
optimization and for accountability.

Measurement overhead should scale with service risk. Best-effort generation
can aggregate telemetry per request. Agentic services require per-step spans
and tool-commit records. Embodied services require per-loop timing and
phase-transition markers, but do not need to expose raw sensory content to the
network. Endpoint-generated value classes and confidential telemetry
aggregation can preserve layering and privacy.

Four rules prevent misleading QoS results.
\begin{enumerate}
\item \emph{Use percentiles and violation runs, not only a mean.} The length
of consecutive deadline misses can be more damaging than the same number of
isolated misses.
\item \emph{Condition results on load and task phase.} A low unloaded TTFT or
loop delay says little about a shared edge at admission saturation.
\item \emph{Separate offered, decoded, and useful goodput.} Tokens delivered
after their deadline or rejected by a version mismatch are not useful
goodput.
\item \emph{Report joint events.} For multimodal fusion, individual
reliability does not reveal the probability that all necessary modalities are
fresh and mutually aligned.
\end{enumerate}

For modalities $m\in{\cal M}_i$ required by task $i$, define synchronized
useful goodput over interval ${\cal T}$ as
\begin{equation}
G_i^{\rm sync}=
\frac{1}{|{\cal T}|}
\sum_{k\in{\cal K}_i}
v_k{\bf 1}\!\left\{
t_k^{\rm rx}\le d_k,\,
\max_{m,n\in{\cal M}_i}|a_m-a_n|\le s_i
\right\},
\label{eq:syncgoodput}
\end{equation}
where $v_k$ is a normalized task-value weight, $a_m$ is information age, and
$s_i$ is the permitted modality skew. This metric does not replace physical
task evaluation, but it is a more faithful online proxy than raw bit or token
rate.

Root-cause attribution can then use the delay budget in
\eqref{eq:loop_delay}. The network should report both the total distribution
and conditional components: radio delay given channel state, compute delay
given selected model and batch, and control loss given task phase. The
conditional view helps transfer a result across deployments. A resource
policy that succeeds only because one simulator has a fast encoder will
otherwise appear universal.

The measurement contract also enables service assurance. Admission uses
predicted distributions and dependency graphs; runtime monitoring compares
observed joint violations with the SLO; adaptation changes compression,
placement, or period; post-mission audit reconstructs the causal chain. This
closed assurance loop is a prerequisite for meaningful service-level
agreements (SLAs) for distributed LAM agents.

\section{From Generative to Embodied Networks for LAMs: Architecture Evolution}
\label{sec:architecture}

\subsection{Architecture for Generative-LAM Networks}
A generative-LAM network is commonly cloud-centric. User equipment performs
tokenization and lightweight preprocessing; an edge gateway provides
admission, caching, privacy filtering, or small-model inference; and a cloud
cluster hosts the full model. Model, tensor, pipeline, and expert parallelism
distribute inference inside or across clusters. Wireless distributed
mixture-of-experts systems additionally make expert routing dependent on
channel and compute heterogeneity \cite{wdmoe}. Pipeline parallelism and
batching trade communication bubbles against accelerator utilization
\cite{pipeline}.

The main control loop is a slow resource-management loop: predict request and
token lengths, place models, reserve accelerators and backhaul, batch
compatible requests, and route overflow. Split inference places early layers
on the device and later layers at the edge, transmitting intermediate
features. Progressive feature transmission allows a receiver to stop once
confidence is sufficient \cite{progressive}. This architecture is effective
when the user transaction is largely independent and deadlines are elastic.
It becomes inefficient when every intermediate action depends on fresh tool
or environment state.

\subsection{Architecture for Agentic-LAM Networks}
Agentic architecture adds a stateful orchestration plane. A planner or
supervisor discovers capabilities, selects models and tools, maintains a task
graph, and commits results to short- and long-term memory. Executors can be
distributed over devices, edge sites, enterprise services, and clouds.
Messages carry not only tokens but role, goal, schema, provenance, confidence,
and dependency. Multi-agent systems use debate, vote, expert routing, or
hierarchical decomposition; the best topology depends on task structure and
the correlation of model errors \cite{agentsmobile}.

This design resembles a distributed transaction system but has probabilistic
components. It requires:
\begin{enumerate}
\item capability and context discovery before routing;
\item session affinity for key--value caches and private memory;
\item checkpoint, retry, and idempotence semantics for tool calls;
\item verification paths for high-impact actions; and
\item joint compute--network--storage admission for the complete task rather
than each RPC.
\end{enumerate}
The architecture may keep a small model on-device for privacy and rapid
screening, escalate uncertain inputs to a remote large model, and invoke
specialists only when their expected information gain exceeds their latency
and cost. Uncertainty-aware hybrid inference is one realization of this
principle \cite{hybrid}.

\subsection{Architecture for Embodied-LAM Networks}
Embodied systems require a hierarchy of loops (Fig.~\ref{fig:architecture}).
A local reflex/control plane executes high-rate stabilization and safety
interlocks. An edge cognition plane performs multimodal fusion, VLA inference,
local mapping, and coordination with deterministic access. A cloud learning
plane provides large-model reasoning, fleet memory, simulation, and
non-real-time model updates. A slow cloud model should not directly occupy a
hard real-time loop; it supplies goals, policies, latent plans, or verified
skill updates to an edge/local controller.

\begin{figure*}[t]
\centering
\setlength{\tabcolsep}{3pt}
\renewcommand{\arraystretch}{1.2}
\begin{tabular}{c}
\fbox{\parbox{0.88\textwidth}{\centering
\textbf{Cloud learning and global intelligence (seconds--hours)}\\
foundation-model training/updates $\mid$ fleet memory $\mid$ simulation and
digital twins $\mid$ global task planning}}\\[-1pt]
$\Updownarrow$ \small model/skill updates, compressed context, audit traces\\[-1pt]
\fbox{\parbox{0.88\textwidth}{\centering
\textbf{Edge cognition and multi-agent coordination (milliseconds--seconds)}\\
multimodal fusion $\mid$ VLA/world-model inference $\mid$ model/expert sharing
$\mid$ task/physical interaction graph $\mid$ deterministic radio scheduling}}\\[-1pt]
$\Updownarrow$ \small synchronized features, state/action tokens, deadlines,
uncertainty\\[-1pt]
\fbox{\parbox{0.88\textwidth}{\centering
\textbf{On-device reflex, control, and safety (sub-ms--milliseconds)}\\
sensor acquisition $\mid$ state estimation $\mid$ servo control $\mid$
collision avoidance $\mid$ fail-safe fallback}}\\[3pt]
\multicolumn{1}{c}{\small
\textit{Cross-layer orchestrator: token value + radio state + compute queues +
control phase}}
\end{tabular}
\caption{Hierarchical architecture for embodied LAM agents. Hard safety and
stabilization remain local; the edge closes latency-critical cognition and
coordination loops; the cloud supplies global reasoning and learning.}
\label{fig:architecture}
\end{figure*}

The architecture is also \emph{partially centralized}. If two robots have
neither task dependence nor physical interference, merging their observations
in one model can add quadratic attention cost without useful information.
When they jointly carry an object or occupy the same collision region,
shared inference can improve consistency and reduce duplicated sensing. Let
${\cal G}_t=({\cal V},{\cal E}_t)$ be a time-varying interaction graph with
\begin{equation}
w_{ij}(t)=\rho_{\rm task} I_{ij}^{\rm task}(t)
+\rho_{\rm phy} I_{ij}^{\rm phy}(t)
+\rho_{\rm info} I_{ij}^{\rm info}(t),
\label{eq:interaction}
\end{equation}
where the three terms quantify task dependency, physical interference, and
complementary information. Clustering or sparsifying ${\cal G}_t$ determines
which agents share a model or exchange tokens. The resulting groups can run
at different loop periods, producing an asynchronous, multi-rate networked
control system.

\begin{table*}[t]
\caption{Architecture Evolution and the New State Introduced at Each Stage}
\label{tab:architecture}
\centering
\footnotesize
\renewcommand{\arraystretch}{1.18}
\setlength{\tabcolsep}{3pt}
\begin{tabular}{|p{0.09\textwidth}|p{0.18\textwidth}|p{0.17\textwidth}|p{0.18\textwidth}|p{0.19\textwidth}|p{0.11\textwidth}|}
\hline
\textbf{Stage} & \textbf{Typical placement} & \textbf{Control/orchestration}
& \textbf{State that must persist} & \textbf{Dominant cross-domain coupling}
& \textbf{Timescale} \\ \hline
Generative &
Device front end; edge gateway/cache; cloud accelerator cluster &
Admission, batching, cache-aware routing, model/expert placement &
Request, model version, KV cache, queue and accelerator state &
Wireless/backhaul rate with memory and compute utilization &
Request and decode-step \\ \hline
Agentic &
Device supervisor; distributed edge/cloud models, tools, and stores &
Task-graph scheduling, capability discovery, verification, retry/rollback &
Goal, plan dependencies, context/memory version, tool side effects,
provenance &
Network--compute--storage--tool transaction &
Step, workflow, and session \\ \hline
Embodied &
Local reflex/safety; edge VLA/world model; cloud fleet learning &
Interaction-graph formation, multi-rate loop scheduling, model sharing,
fail-safe switching &
Physical state, modality timestamps, uncertainty, action history, safety
envelope &
Sensing--communication--computing--control and inter-agent dynamics &
Sub-ms reflex to long-term learning \\ \hline
\end{tabular}
\end{table*}

\subsection{Architecture Design Principles}
Three principles follow. First, \emph{place deadlines before models}: determine
which loop can tolerate cloud, edge, or only local execution before selecting
model size. Second, \emph{centralize only useful mutual information}: shared
models should be activated by task or physical coupling rather than by agent
density alone. Third, \emph{retain a safe degraded mode}: if the edge model,
radio link, or clock synchronization fails, the device must fall back to a
local controller that maintains a safe set, even if task progress stops.

\subsection{Service Lifecycle and Timescale Separation}
Architecture must cover the service lifecycle, not only the inference data
path. Fig.~\ref{fig:lifecycle} separates five stages. During
\emph{registration}, endpoints advertise model/tokenizer versions, sensor and
actuator capabilities, compute envelopes, safety fallbacks, and trust
attestation. During \emph{admission}, the orchestrator maps a requested task
to a service graph, predicts complete resource demand, selects placements,
and rejects or degrades a task whose robust SLO is infeasible. During
\emph{execution}, fast schedulers allocate token, radio, and compute resources
under the admitted envelope. During \emph{recovery}, endpoints switch models,
paths, or controllers and preserve transaction semantics. During
\emph{learning}, traces update workload, token-value, channel, and physical
interaction models.

\begin{figure*}[t]
\centering
\setlength{\tabcolsep}{2.2pt}
\renewcommand{\arraystretch}{1.05}
\begin{tabular}{c@{\,$\rightarrow$\,}c@{\,$\rightarrow$\,}c@{\,$\rightarrow$\,}c@{\,$\rightarrow$\,}c}
\fbox{\parbox[c][0.78in][c]{0.15\textwidth}{\centering\textbf{Register}\\
capability\\version/trust\\fallback}}
&
\fbox{\parbox[c][0.78in][c]{0.15\textwidth}{\centering\textbf{Admit}\\
service graph\\placement\\risk budget}}
&
\fbox{\parbox[c][0.78in][c]{0.15\textwidth}{\centering\textbf{Execute}\\
token value\\radio/compute\\loop monitor}}
&
\fbox{\parbox[c][0.78in][c]{0.15\textwidth}{\centering\textbf{Recover}\\
reroute/retry\\safe mode\\rollback}}
&
\fbox{\parbox[c][0.78in][c]{0.15\textwidth}{\centering\textbf{Learn}\\
causal trace\\model update\\SLO revision}}
\end{tabular}
\caption{Lifecycle of a distributed LAM-agent service. Execution is only one
stage; safe operation depends on versioned registration, end-to-end admission,
explicit recovery, and trace-driven learning.}
\label{fig:lifecycle}
\end{figure*}

The lifecycle contains at least three control timescales.
\begin{itemize}
\item At a \emph{planning timescale} (minutes to hours), models and adapters
are placed, capacity and slices are reserved, and fleet-wide policies are
updated. Forecast error dominates.
\item At a \emph{task timescale} (hundreds of milliseconds to minutes),
agents/tools are selected, model-sharing groups and loop periods change, and
context migrates. Task and physical topology dominate.
\item At a \emph{loop/slot timescale} (sub-millisecond to tens of
milliseconds), features and actions are scheduled, short packets are coded,
and compute slots are aligned. Channel, queue, and release-time uncertainty
dominate.
\end{itemize}
Optimizing every variable at the fastest timescale is neither feasible nor
stable. Slow decisions should pass resource and risk budgets downward; fast
controllers should return feasibility margins and violation statistics.
When fast scheduling cannot maintain an admitted envelope, it should trigger
a defined degradation---smaller model, fewer views, slower noncritical loop,
or local safe mode---rather than silently accumulating queue.

Mobility makes state migration a first-class operation. A generative session
can move after transferring a key--value cache, subject to a short pause. An
agentic session also needs memory and tool-transaction state. An embodied
session cannot pause an unsafe plant while a large cache migrates. Predictive
handover should replicate only the state needed to continue the near-term
loop, keep stabilization local, and transfer large context in the background.
Multi-tier computation offloading \cite{multitier} and ultra-low-latency edge
inference \cite{edgeinference}, together with reliability-aware edge
offloading \cite{offloading}, provide resource-allocation foundations, but
embodied migration adds continuity of action and safety state.

Finally, administrative and physical trust domains need not coincide. A
warehouse owner, robot vendor, network operator, and cloud-model provider may
control different lifecycle stages. The architecture therefore needs signed
capability manifests, explicit data-use policy, model provenance, and
auditable responsibility at every boundary. These controls cannot be added
only at the application layer if network scheduling depends on declared token
value.

\section{Improving Networks for LAMs: Key Enabling Technologies}
\label{sec:enablers}

We group mechanisms into three interacting planes. The \emph{token plane}
decides what representation to compute and communicate. The
\emph{transport plane} decides how to deliver it by a deadline. The
\emph{coordination plane} decides which agents/models should share
information and how network and compute resources close the loop.

\subsection{Enablers for Generative-LAM Networks}
\subsubsection{Compression and token communications}
Quantization, structured pruning, low-rank adaptation, key--value-cache
compression, and top-$k$/top-$p$ decoding reduce memory, computation, or
output entropy. Token communications elevates tokens to semantic
communication units and uses cross-modal context to reconstruct masked or
corrupted content \cite{tokcom}. Compared with conventional source coding, a
receiver equipped with a compatible foundation model can exploit learned
priors. This can reduce payload but transfers burden to shared model state,
tokenizer agreement, uncertainty calibration, and semantic error detection.

For sources $Z_1,\ldots,Z_M$ and task goal $g$, the relevant multi-source
rate--distortion function is conceptually
\begin{equation}
\begin{split}
R_g(D)&=\min_{p(\hat z|z_1,\ldots,z_M)}
I(Z_1,\ldots,Z_M;\hat Z\mid g),\\[-2pt]
&\hspace{8pt}\text{s.t.}\quad
\mathbb{E}[D_{\rm task}(Z,\hat Z\mid g)]\le D.
\end{split}
\label{eq:rd}
\end{equation}
Unlike pixel mean-square error, $D_{\rm task}$ can be nonconvex and
state-dependent. Cross-modal and multiview redundancy are not additive:
removing a camera token may be harmless when lidar observes the same object
but damaging when glare invalidates the lidar-camera correspondence.

\subsubsection{Split, distributed, and speculative inference}
Split inference selects a cut layer and transmits intermediate features.
Tensor parallelism distributes layer operations and repeatedly aggregates
partial results; over-the-air computation can exploit wireless superposition
to accelerate aggregation \cite{distributedllm,aircompneural}. Sparse-feature
over-the-air fusion similarly combines multiview features before
reconstruction or inference \cite{airfusion}. These schemes are attractive
when synchronized devices observe a common scene, but analog aggregation
noise and stragglers must be evaluated at the task output. Intent-aware VLM
split computing for disaster response is an early example of extending the
placement decision toward embodied context \cite{avery}.

Speculative decoding lets a small ``draft'' model propose tokens and a large
model verify them, preserving the target model distribution while reducing
serial decode cost \cite{specdecode}. Wireless implementations must jointly
select the draft model, speculation length, offloading location, and radio
resources. A low draft acceptance rate wastes both inference and
communication. For generative services, the decision can be driven by
expected decode speedup; for embodied services it must also account for action
deviation and deadline risk \cite{actiondeviation}.

\subsection{Enablers for Agentic-LAM Networks}
\subsubsection{Reliable collaboration and model routing}
Multi-model collaboration can improve coverage and verification through
expert routing, debate, vote, or sequential refinement. The network problem
is to choose the smallest collaboration graph that reaches a target task
reliability before a deadline. Correlated models should not be treated as
independent replicas: if they share training data or failure modes, majority
vote offers less diversity than its communication cost suggests.

Adaptive repeat-query systems learn whether another response is likely to
move the latent solution toward a correct concept \cite{garq}. More generally,
an agent should stop when the expected value of another message is lower than
its marginal compute, token, and latency cost. Network schedulers can use
confidence and disagreement as hints, but final verification must remain
application-controlled because confidence is not a safety certificate.

\subsubsection{State, freshness, and transactional transport}
Agent traffic requires protocol support for a task identifier, parent step,
model/tool version, context hash, timeout, and idempotence policy. Memory and
retrieval should be scheduled by \emph{age of relevant information}, not
storage age alone. A newer observation can be irrelevant while an older
constraint remains binding. For irreversible tools, an ``exactly once''
semantic cannot be inferred from transport acknowledgments; the tool and
orchestrator need a commit protocol and audit log.

Compute-aware admission is equally important. Reserving radio resources for a
tool call without an available model replica merely shifts queuing delay.
Conversely, prewarming a model without a feasible wireless path wastes
accelerator memory. Joint admission can operate at two timescales: slow model
placement and slice reservation, followed by fast task-graph scheduling and
rerouting.

\subsection{Enablers for Embodied-LAM Networks}
\subsubsection{Control-loss-aware token selection}
Embodied compression should estimate the marginal increase in control cost
caused by removing, quantizing, delaying, or reusing token $z_k$:
\begin{equation}
s_k(t)=
\frac{\mathbb{E}[\Delta J_{\rm c}\mid z_k,x_t,g_t]}
{\Delta b_k+\alpha\Delta f_k+\beta\Delta e_k},
\label{eq:token_value}
\end{equation}
where the denominator is saved communication, compute, and energy resource.
Tokens with high $s_k$ receive stronger coding or earlier service. The score
changes with control phase: background visual patches can be pruned during
free motion, contact geometry becomes critical at grasp time, and action
reuse is safest inside stable segments. Channel-wise VLA quantization studies
also indicate that channels have unequal action sensitivity \cite{qvla}.

A practical estimator can combine attention/gradient saliency, model
uncertainty, object and contact semantics, and counterfactual rollouts in a
world model. None alone is sufficient. Attention can be poorly calibrated to
causal task value; gradients are local; and rollouts inherit model bias.
Bregman divergences provide a family of representation losses, while
task-conditioned rate--distortion analysis connects them to resource
allocation. The open theoretical problem is to bound physical control loss
when the VLA mapping is nonlinear and the source modalities are dependent.

\subsubsection{Finite-blocklength and deterministic transport}
URLLC techniques provide a foundation for short embodied messages.
For blocklength $n$, signal-to-noise ratio $\gamma$, and error target
$\varepsilon$, the finite-blocklength rate is approximated by
\begin{equation}
R(n,\varepsilon)\approx C(\gamma)
-\sqrt{\frac{V(\gamma)}{n}}Q^{-1}(\varepsilon)
+\frac{\log_2 n}{2n},
\label{eq:fbl}
\end{equation}
where $C$ and $V$ are channel capacity and dispersion
\cite{polyanskiy}. Finite-blocklength cell-free transmission and joint
uplink/downlink allocation further expose hard-deadline and energy tradeoffs
for multiple devices \cite{cellfree,resourceurlcc}. This exposes the
bandwidth--latency--reliability tradeoff hidden by asymptotic capacity.

Interface diversity, multi-connectivity, packet duplication, and optimized
HARQ reduce outage or retransmission delay \cite{interfacediversity,harq}.
Preemptive puncturing allows urgent traffic to overwrite scheduled broadband
resources, while slicing isolates service classes
\cite{puncturing,ranslicing,mobileran}. Two-timescale QoS-aware multiplexing
is useful when slow reservation must anticipate fast allocation under
nonstationary demand \cite{unifiedqos}. Deterministic-delay work further
optimizes jitter, hard-deadline violation, and high-order delay statistics
under finite-blocklength transmission \cite{mixedjitter,deterministic}.

However, direct transplantation to embodied traffic has two limitations.
First, traditional models often assume independent sources and a fixed
utility loss for punctured broadband packets. Multimodal embodied sources are
coupled, so simultaneous degradation can remove all views of a critical
object. Second, radio delay is only one term in
\eqref{eq:loop_delay}. An end-to-end deterministic service must coordinate
sensing release times, encoder and accelerator queues, radio slots, receiver
inference, and actuation.

\subsubsection{Semantic and goal-oriented communication}
Goal-oriented communication schedules information by its effect on a task.
Value and age of information have been applied to robotic waypoint delivery
\cite{waypoint}. Integrated sensing and edge AI similarly co-designs sensing,
communication, and inference around task performance \cite{isea}. For
embodied LAMs, semantic communication should protect control objectives,
cross-modal complementarity, and multiview complementarity. It should also
report uncertainty: hallucinated recovery of a missing object is unacceptable
when a safety controller interprets it as observation. Rate-limited
closed-loop semantic communication offers a direct bridge from task
representations to distributed sensing and control \cite{semanticclosedloop}.

A useful design is \emph{loss-controlled non-orthogonal reuse}. Two sources
may overlap in radio resources if their joint task distortion remains below a
bound; sources covering the same hazard should instead receive diversity.
This reverses a common heuristic: highly correlated data can be aggressively
compressed for efficiency, but correlation alone does not justify common-mode
failure.

\subsubsection{Communication--control co-design and dynamic model sharing}
Wireless networked control studies optimize bandwidth, power, uplink/downlink
ratio, sampling period, and control policy jointly. Recent work considers
sensing--communication--computing--control loops for robotic digital twins
\cite{fang}, communication--control codesign for large-scale systems
\cite{pang}, and communication allocation for system identification
\cite{systemid}. Energy-aware wireless communication--control co-design
further shows that loop performance and radio efficiency should be optimized
together \cite{energywncs}. These results show that optimizing physical
control or identification can produce allocations very different from
throughput maximization.

Most tractable theory assumes linear dynamics and a common period. Embodied
LAM systems add black-box nonlinear policies, partial observability,
asynchronous periods, and time-varying inter-agent coupling. A robust
orchestrator must jointly choose model-sharing groups ${\cal C}_t$, periods
$h_i$, token sets ${\cal K}_{i,t}$, bandwidth $b_{i,t}$, and compute
$f_{i,t}$:
\begin{subequations}
\begin{align}
\min\quad &\sum_i \omega_i J_{{\rm c},i}
+\lambda_B\sum_i b_i+\lambda_F\sum_i f_i+\lambda_E E,\\
\text{s.t.}\quad&
\Pr\{T_{i,t}^{\rm loop}\le h_i,\,
|\Delta T_{i,t}|\le\epsilon_i h_i\}\ge1-\delta_i,\\
&D_{{\rm task},i}({\cal K}_{i,t},{\cal C}_t)\le D_i^{\max},\\
&\sum_i b_{i,t}\le B,\quad \sum_i f_{i,t}\le F,\\
&{\cal C}_t\in{\cal C}_{\rm safe}.
\end{align}
\label{eq:joint_problem}
\end{subequations}
Slow graph clustering and model placement can be separated from fast token,
radio, and compute scheduling, but the two timescales must exchange
feasibility margins. Distributionally robust or risk-sensitive optimization
is promising because empirical delay and control-loss distributions shift
with scenes and tasks.

\begin{table*}[t]
\caption{Enabling Technologies, Benefits, and Transfer Gaps}
\label{tab:enablers}
\centering
\footnotesize
\renewcommand{\arraystretch}{1.17}
\setlength{\tabcolsep}{3.1pt}
\begin{tabular}{|p{0.14\textwidth}|p{0.17\textwidth}|p{0.19\textwidth}|p{0.19\textwidth}|p{0.21\textwidth}|}
\hline
\textbf{Technology} & \textbf{Primary decision} & \textbf{Benefit}
& \textbf{Representative objective} & \textbf{Gap for embodied multi-agent
LAMs} \\ \hline
Token pruning, quantization, reuse &
Which tokens/channels/steps to retain and at what precision &
Lower payload, memory, and inference cost &
Reconstruction, generation quality, success rate &
Need state- and action-conditioned control-loss bounds; small actions can be
stability-critical \\ \hline
Split/distributed inference and AirComp &
Cut, placement, participating devices, aggregation/transceiver &
Parallel compute and reduced raw-data transfer &
Feature error, latency, energy, inference accuracy &
Stragglers and analog error must be related to asynchronous closed-loop cost
\\ \hline
Speculative/hybrid inference &
Draft model, speculation length, local/remote escalation &
Fewer serial large-model decoding steps &
Acceptance rate, expected latency, generation fidelity &
Action deviation and deadline-tail risk replace text-distribution fidelity
as the decisive constraint \\ \hline
Finite-blocklength, diversity, HARQ &
Blocklength, coding, replicas, interfaces, retransmission &
Reliable delivery under short deadlines &
Packet error, energy, spectrum, one-way latency &
Must include sensing, compute, and actuation tails and token-specific
deadlines \\ \hline
Puncturing, slicing, deterministic scheduling &
Priority, reserved bandwidth, preemption, slot/queue schedule &
Isolation and bounded delay/jitter for mixed traffic &
Deadline violation and broadband utility loss &
Utility loss is coupled across modalities, views, and agents rather than
i.i.d. per packet \\ \hline
Semantic/goal-oriented communication &
What task-relevant representation and update to transmit &
Lower rate for a target inference/control quality &
Age/value of information, task distortion &
Generative recovery can hallucinate; uncertainty and common-mode safety loss
must be bounded \\ \hline
Communication--control co-design and model sharing &
Loop period, resources, controller/model group, collaboration graph &
Direct optimization of physical task and fleet cost &
LQR cost, identification error, long-term control utility &
Nonlinear black-box policies, multi-rate loops, and time-varying
task/physical coupling remain open \\ \hline
\end{tabular}
\end{table*}

\subsection{Cross-Layer Synthesis: From ``Compress Then Protect'' to ``Value Then Coordinate''}
The technologies in Table~\ref{tab:enablers} should not be composed as
independent modules. Selecting tokens changes packet size and inference time;
transport decisions change observation age; model sharing changes both source
correlation and compute topology; all three change physical control cost.
The appropriate workflow is:
\begin{enumerate}
\item infer the task phase, safety envelope, and interaction graph;
\item estimate marginal token/action value and uncertainty;
\item choose local, shared-edge, or cloud inference and a safe fallback;
\item jointly schedule sensing, compute, and short-packet transmission; and
\item observe task outcome and update the value and delay models.
\end{enumerate}
This ``value then coordinate'' loop is the network counterpart of embodied
intelligence: resources are allocated according to how information changes
action, not merely according to how many bytes it occupies.

\section{Case Study: Multi-Agent Robotic Rescue}
\label{sec:case}

\subsection{Scenario and Service Graph}
Consider a post-disaster site in which three robots must remove a beam that
blocks access to an injured person, avoid a fragile temporary support, and
then assist with first aid. The scenario, adapted from the project material
underlying this survey, deliberately combines high-data-rate perception,
physical coupling, and safety-critical manipulation. Each robot carries RGB
and depth cameras, lidar, inertial sensing, and joint encoders. Local servo
controllers maintain balance; an edge VLA/world model supplies manipulation
actions; a cloud model supplies non-real-time semantic reasoning and fleet
knowledge.

The task has four phases:
\begin{enumerate}
\item \emph{search and mapping}: robots explore largely independently and
exchange sparse map/object tokens;
\item \emph{beam assessment}: multiview perception is fused to identify load,
grasp points, and the temporary support;
\item \emph{cooperative lift}: task and physical edges in
${\cal G}_t$ become dense, activating shared edge inference and synchronized
action delivery;
\item \emph{first aid}: fine manipulation receives a strict loop SLO while
background mapping is demoted.
\end{enumerate}
This is not a claim that one fixed delay target fits all robots. Rather, it
shows why service classes must change within one mission.

\begin{table*}[t]
\caption{Task-Phase-Aware Orchestration in the Rescue Case Study}
\label{tab:case}
\centering
\footnotesize
\renewcommand{\arraystretch}{1.2}
\setlength{\tabcolsep}{3.1pt}
\begin{tabular}{|p{0.13\textwidth}|p{0.18\textwidth}|p{0.18\textwidth}|p{0.18\textwidth}|p{0.23\textwidth}|}
\hline
\textbf{Phase} & \textbf{Dominant information} & \textbf{Coordination mode}
& \textbf{Differentiated SLO} & \textbf{Network action} \\ \hline
Search/mapping &
RGB/depth/lidar features, local pose, low-rate discoveries &
Independent local control; sparse peer map exchange &
Perception coverage and freshness; elastic cloud reasoning &
Prune background visual tokens; cache maps at edge; opportunistic model/tool
calls \\ \hline
Beam assessment &
Multiview geometry, support classification, grasp candidates, uncertainty &
Edge fusion among robots with complementary views &
Bound modality skew and common-object semantic distortion &
Allocate view diversity; avoid jointly dropping tokens that observe the same
hazard; verify recovered semantics \\ \hline
Cooperative lift &
Force/torque, joint state, beam pose, synchronized action/intent &
Temporary shared VLA/world model plus local safety controllers &
High loop-deadline reliability, bounded jitter and inter-robot skew &
Reserve deterministic slice; duplicate critical state/action packets; align
compute and radio slots; preempt background traffic \\ \hline
First aid &
Close-range vision, tool pose, contact state, fine actions &
One precision group; other robots act as sensing/safety peers &
Low action error and safety violation; moderate semantic planning latency &
Prioritize contact-transition tokens; disable unsafe action reuse; preserve
local fail-safe and audit trail \\ \hline
\end{tabular}
\end{table*}

\subsection{Cross-Layer Orchestration}
At mission start, the orchestrator builds ${\cal G}_t$ from task dependencies
and predicted collision/interaction regions. Sparse components run separate
models; the cooperative-lift component shares an edge model because each
robot's action changes the others' observations and feasible actions. The
model-sharing decision is recomputed at a slower timescale than radio
scheduling to avoid oscillation.

Within a group, token scores in \eqref{eq:token_value} separate three classes.
\emph{Class A} contains safety constraints, contact transitions, force/torque
events, and final actions; it receives deterministic scheduling and, when
needed, path diversity. \emph{Class B} contains task-relevant visual and
geometric features; it receives loss-controlled compression and synchronized
fusion. \emph{Class C} contains background context, rationales, logs, and
model updates; it uses best effort and can be deferred. Classification is
revisited after every task-phase transition.

The scheduler solves a practical approximation to
\eqref{eq:joint_problem}. It first reserves local compute and radio capacity
for Class A under the robust-loop constraint, then maximizes the marginal
task value of Class B, and finally serves Class C. If the edge queue threatens
the loop budget, the system shrinks the shared group, selects a smaller VLA,
reduces visual-token density, or hands control to the local safe controller.
Simply increasing transmit power cannot fix an overloaded inference queue.

\subsection{Evaluation Protocol and Reproducibility}
The project material reports an Isaac Sim--GR00T testbed connected through an
RPC channel that can inject bandwidth limits, delay, jitter, and errors. A
rigorous evaluation should publish the model checkpoint, tokenizer, scene and
random seeds, sensor rates, compute hardware, channel traces, queueing
discipline, clock synchronization, and controller fallback. Each network
configuration should be evaluated over complete task episodes, not isolated
frames.

Four ablations identify causal gain:
\begin{enumerate}
\item byte-aware versus task/token-aware scheduling;
\item independent versus task/physical-graph-based model sharing;
\item mean-delay optimization versus robust delay--jitter constraints; and
\item raw/full transmission versus spatial perception pruning plus
phase-aware action reuse.
\end{enumerate}
Report task success with confidence intervals, collision/safety violations,
trajectory and action error, $J_{\rm c}$, loop-delay percentiles, jitter,
modality skew, bandwidth, energy, and accelerator utilization. Results must
also include \emph{failure slices}: occlusion, a misleading view, burst
wireless loss, edge straggler, clock drift, and an abrupt phase transition.
This protocol prevents a bandwidth gain obtained by silently increasing
physical risk from being presented as a network improvement.

\section{Challenges and Future Directions}
\label{sec:future}

\subsection{Token Recognition and Sorting}
Today's network typically observes encrypted packets, endpoints, and coarse
application identifiers, not the task value of individual tokens. Yet the
survey's proposed differentiation requires modality, deadline, dependency,
uncertainty, and control phase. Deep packet inspection is neither scalable nor
compatible with end-to-end confidentiality.

A promising solution is an explicit, privacy-preserving token descriptor
generated by the trusted endpoint. It can carry a compact class rather than
content: service stage, modality, generation/measurement timestamp, deadline,
dependency group, reliability tier, and uncertainty/value bin. The descriptor
must be authenticated because an untrusted application could mark every token
as safety-critical. Admission control should police declared value against a
contract and observed traffic.

Sorting is also a causal inference problem. Attention, activation magnitude,
entropy, or gradient saliency can correlate with task value but fail under
distribution shift. Future work should develop counterfactual estimators:
what change in task or control cost results if this token is delayed,
quantized, predicted, or omitted? World-model rollouts can approximate the
answer, while conformal or distributionally robust bounds can express
uncertainty. The difficult case is a rare but safety-critical token whose
average value is low. Risk-sensitive sorting must preserve it.

Three granularity levels are needed. Flow-level recognition supports backward
compatibility but is coarse. Chunk-level recognition groups tokens with a
common deadline and coding policy. Token-level recognition offers maximum
efficiency but can introduce excessive header, scheduling, and privacy cost.
An adaptive hierarchy should use token-level metadata only around task
transitions or hazards.

\subsection{Tokenizer Compatibility}
Token communications assumes that sender and receiver agree on token
boundaries, codebooks, embeddings, and model semantics. In practice,
tokenizers differ by model, version, language, modality, and embodiment.
Even identical token IDs are meaningless without a versioned vocabulary and
embedding interpretation. Visual and action tokenizers are especially
unstable because codebooks can change during fine-tuning.

Compatibility has four layers:
\begin{enumerate}
\item \emph{syntactic}: vocabulary and token-ID mapping;
\item \emph{representational}: embedding dimension, quantization, and feature
normalization;
\item \emph{semantic}: equivalent meaning or task effect across models; and
\item \emph{behavioral}: equivalent downstream action within a safety
tolerance.
\end{enumerate}
Only the first two can be solved by a schema registry. Semantic and behavioral
compatibility require calibration data, adapters, or a shared latent
interface. A universal tokenizer may simplify networking but can sacrifice
model efficiency and lock the ecosystem to one representation.

Future protocols should negotiate a tokenizer/model manifest containing a
cryptographic version hash, modality, vocabulary or codebook identifier,
embedding format, supported adapters, calibration range, and fallback codec.
Gateways can translate tokens through learned adapters, but translation adds
latency and an error that must be included in
\eqref{eq:task_distortion}. For safety-critical actions, a receiver should
reject an unknown or out-of-range codebook rather than guess.

Model updates create a consistency problem similar to a rolling protocol
upgrade. During a fleet update, old and new tokenizers may coexist. Dual
encoding is robust but expensive; atomic migration is operationally hard.
Version-aware routing and shadow validation can support a staged transition.
A major research need is a formal bound connecting embedding alignment error
to task and control degradation.

\subsection{Token-Native Protocols}
IP networks are byte- and packet-native; RPC frameworks are
request-native. LAM services need a transaction and token/chunk abstraction
without discarding the scalability and isolation of existing layers. A
token-native protocol should expose only the information necessary for
scheduling and recovery, and should remain usable when intermediate networks
are unaware of tokens.

The protocol data unit could include task/session and step/loop identifiers,
model and tokenizer versions, modality, timestamp and deadline, dependency,
value or reliability class, and uncertainty.
It should support partial reliability, deadline-aware abandonment,
application-level forward error correction, semantic acknowledgments, and
multi-path duplication. A semantic acknowledgment confirms that a receiver
obtained sufficient task information, not simply all bytes. It must not allow
a generative model to acknowledge a hallucinated reconstruction as a verified
sensor observation.

Congestion control also changes. Retransmitting an expired action wastes
capacity and can be unsafe; dropping a low-value context token may be optimal;
and increasing the rate of observation tokens can overload the shared
inference queue. Congestion feedback should therefore combine radio
congestion, compute queue, cache pressure, and deadline miss risk. Scheduling
must prevent value inflation from starving conventional traffic.

Security is inseparable from protocol design. Token metadata leaks user
intent, robot phase, and model identity. Attackers can manipulate value tags,
inject stale actions, exploit tokenizer mismatch, or poison shared context.
Required defenses include authenticated timestamps and manifests, replay
protection, provenance, least-privilege tool tokens, confidential inference,
and safety envelopes checked outside the generative model.

\subsection{Benchmarking, Theory, and Standardization Gaps}
Current benchmarks usually isolate model accuracy, network traces, or control
dynamics. The community needs co-simulation and hardware testbeds that combine
real model serving, wireless queues, clock drift, sensor release patterns, and
physical episodes. Benchmarks should span generative, agentic, and embodied
traffic on one infrastructure so that resource competition and backward
compatibility can be studied. Shared datasets should include task/physical
interaction graphs and time-aligned token-importance labels, while protecting
private sensory content.

Theory must move beyond independent sources, stationary delay, linear plants,
and known dynamics. Promising tools include multi-terminal
task-rate--distortion theory, risk-sensitive and distributionally robust
optimization, event-triggered control, graph sparsification, and online
learning with safety shields. The central missing bound connects four errors:
token representation distortion, wireless delivery uncertainty, model
prediction error, and long-horizon physical control loss.

Standardization should proceed incrementally. Near-term systems can define
LAM service descriptors and telemetry over existing QoS flows. A second step
can standardize capability/tokenizer manifests and compute-aware admission.
Only after interoperability and security experience should fine-grained
token-native transport be exposed across domains. The active 3GPP 6G use-case
study \cite{3gpp22870} provides a venue, but task semantics and safety policy
must remain under application and vertical-domain governance.

\subsection{A Staged Research Roadmap}
The open problems should be attacked in an order that produces falsifiable
interfaces rather than an all-at-once ``AI-native network.'' In the near term,
research should establish \emph{measurement compatibility}. Model-serving,
agent, radio, compute, and robotics platforms need the causal identifiers and
outcome records described in Section~\ref{sec:qos}. A useful milestone is a
public workload in which the same task episode can be replayed under multiple
tokenizers, model placements, channel traces, and queueing policies, with
every deadline miss attributable to one component. At this stage, endpoints
can map rich private state to a small set of conventional QoS flows; no new
over-the-air token format is required.

The next stage is \emph{service and representation interoperability}.
Capability manifests, version negotiation, chunk-level deadline/value
classes, compute-aware admission, and semantic acknowledgments can be
implemented above existing transport. Success should be judged by
cross-vendor portability: an agent should discover a remote model or tool,
verify its schema and tokenizer adapter, predict an end-to-end SLO, and fall
back without corrupting task state. For embodied systems, the corresponding
test is a rolling edge/model update in which local safety remains continuous
and the task either progresses or enters a documented safe stop.

A third stage can introduce \emph{token- and loop-native optimization} inside
the network. Schedulers may use authenticated value/deadline bins, joint
radio--compute congestion feedback, modality-skew constraints, and
task-conditioned partial reliability. The key scientific milestone is not a
higher token rate; it is a reproducible Pareto improvement in physical task
utility per unit cost under distribution shift. Evaluation must include a
strong byte-aware baseline with the same total resources, and must show that
value signaling cannot be gamed to starve other tenants.

Longer term, networks can support \emph{distributed embodied intelligence as a
managed closed-loop service}. Interaction graphs, model-sharing groups, loop
periods, and risk budgets would adapt jointly across a fleet. Achieving this
requires theory that separates predictable learned dynamics from residual
risk, plus certification interfaces between probabilistic models and
deterministic safety controllers. The network should never be asked to infer
an application's safety policy from raw tokens; it should enforce a
machine-verifiable contract supplied by the responsible vertical.

Across all stages, three negative results are as valuable as positive ones.
Researchers should identify when semantic compression loses to conventional
coding after model/adapter overhead, when shared inference loses to independent
control after synchronization and attention cost, and when a tighter radio
tail produces no control gain because compute or sensing dominates. These
boundary conditions turn the taxonomy in this paper into engineering
guidance, prevent token-aware networking from becoming a universal but
unfalsifiable claim, and direct standardization toward mechanisms with
measurable end-to-end value.

\section{Conclusions}
\label{sec:conclusion}
The network evolution from generative to agentic and embodied LAMs is not a
simple increase in traffic volume. It is a change in the unit of service:
from a response, to a stateful task, to a physical closed loop. Accordingly,
QoS progresses from content quality and token latency, through long-horizon
task reliability, to control cost, tail delay, jitter, synchronization, and
safety. Architectures progress from cloud-centric serving, through
distributed transaction orchestration, to hierarchical local--edge--cloud
control with dynamic model sharing.

Existing foundations---efficient inference, semantic communications, URLLC,
and wireless networked control---are necessary but individually insufficient.
Embodied networks require control-loss-aware token selection,
source-coupled deterministic transport, and task/physical-graph-driven
coordination. The practical design rule is to evaluate information by the
action it changes, coordinate radio and compute around the complete loop, and
retain a local safe fallback. Token recognition, tokenizer interoperability,
and token-native protocols are therefore not cosmetic optimizations; they are
the interfaces through which future networks can connect distributed
intelligence to the physical world.

\begin{thebibliography}{50}

\bibitem{3gpp22870}
3GPP, ``Study on 6G use cases and service requirements,'' 3rd Generation
Partnership Project (3GPP), Tech. Rep. 22.870, Rel. 20, 2026.

\bibitem{mmlu}
D. Hendrycks, C. Burns, S. Basart, A. Zou, M. Mazeika, D. Song, and
J. Steinhardt, ``Measuring massive multitask language understanding,'' in
\emph{Proc. ICLR}, 2021.

\bibitem{bigbench}
BIG-bench Authors, ``Beyond the imitation game: Quantifying and extrapolating
the capabilities of language models,'' \emph{Trans. Mach. Learn. Res.}, 2023.

\bibitem{agentbench}
X. Liu \etal, ``AgentBench: Evaluating LLMs as agents,'' in \emph{Proc. ICLR},
2024.

\bibitem{metaworld}
T. Yu \etal, ``Meta-World: A benchmark and evaluation for multi-task and
meta reinforcement learning,'' in \emph{Proc. Conf. Robot Learn.}, 2020,
pp. 1094--1100.

\bibitem{groot}
NVIDIA \etal, ``GR00T N1: An open foundation model for generalist humanoid
robots,'' arXiv:2503.14734, Mar. 2025.

\bibitem{flashvla}
X. Tan, Y. Yang, P. Ye, J. Zheng, B. Bai, X. Wang, J. Hao, and T. Chen,
``Think twice, act once: Token-aware compression and action reuse for
efficient inference in vision-language-action models,'' arXiv:2505.21200,
May 2025.

\bibitem{beyondcloud}
X. Zhang \etal, ``Beyond the cloud: Edge inference for generative large
language models in wireless networks,'' \emph{IEEE Trans. Wireless Commun.},
vol. 24, no. 1, pp. 643--658, Jan. 2025.

\bibitem{garq}
T. Yu, K. Huang, J. Li, H. Pervaiz, G. Zheng, and S. Zhang, ``Edge
heterogeneous LLMs collaborative reliable inference based on adaptive repeat
query,'' in \emph{Proc. IEEE ICC Workshops}, 2026, accepted.

\bibitem{netmcp}
E. Li, H. Du, and K. Huang, ``NetMCP: Network-aware model context protocol
platform for LLM capability extension,'' arXiv:2510.13467, Oct. 2025.

\bibitem{wdmoe}
N. Xue \etal, ``WDMoE: Wireless distributed mixture of experts for large
language models,'' \emph{IEEE Trans. Wireless Commun.}, vol. 25,
pp. 559--572, 2026.

\bibitem{pipeline}
J. Jiang, Z. Chen, H. Shin, and A. Nallanathan, ``Edge inference for large
language models with pipeline parallelism and batching,'' \emph{IEEE Trans.
Commun.}, early access, 2026.

\bibitem{progressive}
Q. Lan, Q. Zeng, P. Popovski, D. G\"und\"uz, and K. Huang, ``Progressive
feature transmission for split classification at the wireless edge,''
\emph{IEEE Trans. Wireless Commun.}, vol. 22, no. 6, pp. 3837--3852,
Jun. 2023.

\bibitem{hybrid}
S. Oh, J. Kim, J. Park, S.-W. Ko, T. Q. S. Quek, and S.-L. Kim,
``Uncertainty-aware hybrid inference with on-device small and remote large
language models,'' in \emph{Proc. IEEE ICMLCN}, 2025, pp. 1--7.

\bibitem{tokcom}
L. Qiao, M. B. Mashhadi, Z. Gao, R. Tafazolli, M. Bennis, and D. Niyato,
``Token communications: A large model-driven framework for cross-modal
context-aware semantic communications,'' \emph{IEEE Wireless Commun.},
vol. 32, no. 5, pp. 80--88, Oct. 2025.

\bibitem{distributedllm}
K. Zhang, H. He, S. Song, J. Zhang, and K. B. Letaief,
``Communication-efficient distributed on-device LLM inference over wireless
networks,'' \emph{IEEE J. Sel. Topics Signal Process.}, vol. 19, no. 7,
pp. 1301--1317, Oct. 2025.

\bibitem{airfusion}
Z. Liu, Q. Lan, and K. Huang, ``Over-the-air fusion of sparse spatial
features for integrated sensing and edge AI over broadband channels,''
\emph{IEEE Trans. Wireless Commun.}, vol. 24, no. 4, pp. 2999--3013,
Apr. 2025.

\bibitem{specdecode}
Y. Leviathan, M. Kalman, and Y. Matias, ``Fast inference from transformers
via speculative decoding,'' in \emph{Proc. ICML}, 2023, pp. 19274--19286.

\bibitem{actiondeviation}
J. Park, Y. Lim, S. Oh, J. Park, J. Choi, and S.-L. Kim, ``Action
deviation-aware inference for low-latency wireless robots,''
arXiv:2510.02851, Nov. 2025.

\bibitem{qvla}
Y. Xu \etal, ``QVLA: Not all channels are equal in vision-language-action
model's quantization,'' arXiv:2602.03782, Feb. 2026.

\bibitem{polyanskiy}
Y. Polyanskiy, H. V. Poor, and S. Verd\'u, ``Channel coding rate in the finite
blocklength regime,'' \emph{IEEE Trans. Inf. Theory}, vol. 56, no. 5,
pp. 2307--2359, May 2010.

\bibitem{interfacediversity}
J. J. Nielsen, R. Liu, and P. Popovski, ``Ultra-reliable low latency
communication using interface diversity,'' \emph{IEEE Trans. Commun.},
vol. 66, no. 3, pp. 1322--1334, Mar. 2018.

\bibitem{harq}
A. Anand and G. de Veciana, ``Resource allocation and HARQ optimization for
URLLC traffic in 5G wireless networks,'' \emph{IEEE J. Sel. Areas Commun.},
vol. 36, no. 11, pp. 2411--2421, Nov. 2018.

\bibitem{puncturing}
A. Anand, G. de Veciana, and S. Shakkottai, ``Joint scheduling of URLLC and
eMBB traffic in 5G wireless networks,'' \emph{IEEE/ACM Trans. Netw.},
vol. 28, no. 2, pp. 477--490, Apr. 2020.

\bibitem{ranslicing}
P. Yang, X. Xi, T. Q. S. Quek, J. Chen, X. Cao, and D. Wu, ``RAN slicing
for massive IoT and bursty URLLC service multiplexing: Analysis and
optimization,'' \emph{IEEE Internet Things J.}, vol. 8, no. 18,
pp. 14258--14275, Sep. 2021.

\bibitem{mixedjitter}
Z. Liu, K. Li, Y. Wang, and P. Zhu, ``Mixed traffic scheduling with latency
and jitter analysis in URLLC industrial automation,'' \emph{IEEE Internet
Things J.}, vol. 12, no. 12, pp. 18820--18835, Jun. 2025.

\bibitem{deterministic}
H. Ding, S. Xu, Z. Xu, R. Xu, and B. Ai, ``Achieving deterministic delay in
industrial Internet of Things: Resource allocation schemes in the finite
blocklength regime,'' \emph{IEEE Trans. Veh. Technol.}, vol. 74, no. 9,
pp. 14192--14206, Sep. 2025.

\bibitem{waypoint}
W. Wu, Y. Yang, Y. Deng, and A. H. Aghvami, ``Goal-oriented semantic
communications for robotic waypoint transmission: The value and age of
information approach,'' \emph{IEEE Trans. Wireless Commun.}, vol. 23, no. 12,
pp. 18903--18915, Dec. 2024.

\bibitem{isea}
Z. Liu \etal, ``Integrated sensing and edge AI: Realizing intelligent
perception in 6G,'' \emph{IEEE Commun. Surveys Tuts.}, early access, 2025.

\bibitem{fang}
X. Fang \etal, ``Sensing-communication-computing-control closed-loop
optimization for 6G digital twin-empowered robotic systems,'' \emph{IEEE J.
Sel. Areas Commun.}, vol. 43, no. 10, pp. 3330--3346, Oct. 2025.

\bibitem{pang}
G. Pang, W. Liu, D. Niyato, B. Vucetic, and Y. Li,
``Communication-control codesign for large-scale wireless networked control
systems,'' \emph{IEEE J. Sel. Areas Commun.}, vol. 43, no. 10,
pp. 3295--3312, Oct. 2025.

\bibitem{systemid}
Z. Han, X. Li, Z. Zhou, K. Huang, Y. Gong, and Q. Zhang, ``Wireless
communication and control co-design for system identification,'' \emph{IEEE
Trans. Wireless Commun.}, vol. 23, no. 5, pp. 4114--4126, May 2024.

\bibitem{unifiedqos}
J. Li \etal, ``A unified QoS-aware multiplexing framework for next-generation
immersive communication with legacy wireless applications,'' \emph{IEEE
Internet Things J.}, vol. 12, no. 14, pp. 28582--28597, Jul. 2025.

\bibitem{cellfree}
Z. Zhang \etal, ``Performance of multidevice downlink cell-free system under
finite blocklength for URLLC with hard deadlines,'' \emph{IEEE J. Sel. Areas
Commun.}, vol. 41, no. 7, pp. 2090--2106, Jul. 2023.

\bibitem{offloading}
C.-F. Liu, M. Bennis, M. Debbah, and H. V. Poor, ``Dynamic task offloading
and resource allocation for ultra-reliable low-latency edge computing,''
\emph{IEEE Trans. Commun.}, vol. 67, no. 6, pp. 4132--4150, Jun. 2019.

\bibitem{multitier}
D. V. Huynh, V.-D. Nguyen, S. Chatzinotas, S. R. Khosravirad, H. V. Poor,
and T. Q. Duong, ``Joint communication and computation offloading for
ultra-reliable and low-latency with multi-tier computing,'' \emph{IEEE J.
Sel. Areas Commun.}, vol. 41, no. 2, pp. 521--537, Feb. 2023.

\bibitem{edgeinference}
Z. Wang, A. E. Kal{\o}r, Y. Zhou, P. Popovski, and K. Huang,
``Ultra-low-latency edge inference for distributed sensing,'' \emph{IEEE
Trans. Wireless Commun.}, early access, 2025.

\bibitem{semanticclosedloop}
G. Pan, A. \"Oz\c{c}elikkale, C. H\"ager, M. F. Keskin, and H. Wymeersch,
``Semantic communication for rate-limited closed-loop distributed
communication-sensing-control systems,'' arXiv:2512.19177, Dec. 2025.

\bibitem{energywncs}
Z. Han, X. Li, G. Zhu, K. Huang, Y. Gong, and Q. Zhang, ``An energy-efficient
wireless communication and control co-design for WNCS,'' \emph{IEEE Trans.
Wireless Commun.}, vol. 25, pp. 1798--1812, 2026.

\bibitem{avery}
R. Bhattacharjya \etal, ``AVERY: Intent-driven adaptive VLM split computing
via embodied self-awareness for efficient disaster response systems,''
arXiv:2511.18151, Mar. 2026.

\bibitem{dadu}
Y. Huang \etal, ``DaDu-Corki: Algorithm-architecture co-design for embodied
AI-powered robotic manipulation,'' in \emph{Proc. ACM/IEEE ISCA}, 2025,
pp. 327--343.

\bibitem{mobileran}
M. Alsenwi, N. H. Tran, M. Bennis, S. R. Pandey, A. K. Bairagi, and
C. S. Hong, ``Intelligent resource slicing for eMBB and URLLC coexistence in
5G and beyond: A deep reinforcement learning based approach,'' \emph{IEEE
Trans. Wireless Commun.}, vol. 20, no. 7, pp. 4585--4600, Jul. 2021.

\bibitem{semanticclassic}
H. Xie, Z. Qin, G. Y. Li, and B.-H. Juang, ``Deep learning enabled semantic
communication systems,'' \emph{IEEE Trans. Signal Process.}, vol. 69,
pp. 2663--2675, 2021.

\bibitem{aircompneural}
Y. Cang, M. Chen, and K. Huang, ``Over-the-air computation for realizing
neural link in in-network AI architectures,'' \emph{IEEE Trans. Wireless
Commun.}, vol. 25, pp. 2373--2388, 2026.

\bibitem{resourceurlcc}
K. Li, P. Zhu, Y. Wang, F.-C. Zheng, and X. You, ``Joint uplink and downlink
resource allocation toward energy-efficient transmission for URLLC,''
\emph{IEEE J. Sel. Areas Commun.}, vol. 41, no. 7, pp. 2176--2192,
Jul. 2023.

\bibitem{agentsmobile}
X. Zheng \etal, ``Exploring multi-agent dynamics for generative AI and large
language models in mobile edge networks,'' \emph{IEEE Wireless Commun.},
vol. 32, no. 6, pp. 69--78, Dec. 2025.

\bibitem{modeloffload}
Q. Chen, C. Yang, S. Lan, L. Zhu, and Y. Zhang, ``Two-stage evolutionary
search for efficient task offloading in edge computing power networks,''
\emph{IEEE Internet Things J.}, vol. 11, no. 19, pp. 30787--30799,
Oct. 2024.

\end{thebibliography}

\end{document}
